\documentclass[10pt, a4paper]{article}
\usepackage{fontspec} 

% DOCUMENT LAYOUT
\usepackage{geometry} 
\geometry{a4paper, textwidth=5.5in, textheight=8.5in, marginparsep=7pt, marginparwidth=.6in}
\setlength\parindent{0in}
\usepackage{adforn}
\newcommand{\FEUP}{Faculty of Engineering, University of Porto}

% Identifying information
\newcommand{\Title}{Curriculum Vit\ae}
\newcommand{\FirstName}{João Pedro}
\newcommand{\LastName}{Dias}
\newcommand{\Initials}{J.P.}
\newcommand{\MyName}{\FirstName\ \LastName}
\newcommand{\Me}{\textbf{\LastName, \Initials}}  % For citations
\newcommand{\Email}{jpmdias@fe.up.pt}
\newcommand{\Website}{jpdias.me}
\newcommand{\ORCID}{0000-0001-9066-6436}
\newcommand{\Affiliation}{
    \adforn{10} Kuehne+Nagel\\
    \adforn{10} Department of Informatics Engineering
    \\
    \FEUP
}
\newcommand{\Address}{Porto, Portugal}

\usepackage{marginnote}
\newcommand{\amper{}}{\chardef\amper="E0BD }
\newcommand{\years}[1]{\marginnote{\raggedleft\footnotesize #1}}
\renewcommand*{\raggedleftmarginnote}{}
\setlength{\marginparsep}{10pt}
\reversemarginpar
\usepackage[leftmargin=0pt,labelsep=10pt,itemsep=2pt,parsep=2pt]{etaremune}

\usepackage[p,osf]{cochineal}
\usepackage[type1]{cabin}
\usepackage[cochineal,bigdelims,cmintegrals,vvarbb]{newtxmath}
\usepackage[zerostyle=c,scaled=.94]{newtxtt}
\usepackage[cal=boondoxo]{mathalfa}

% Control the font size
\usepackage{anyfontsize}
\usepackage[maxbibnames=99,maxcitenames=99]{biblatex}
%\bibliographystyle{plain}
\bibliography{reference-personal}

% Template variables for styling
\newcommand{\TablePad}{\vspace{-0.4cm}}
\newcommand{\SoftwareTitle}[1]{{\bfseries #1}}
\newcommand{\TableTitle}[1]{{\fontsize{12pt}{0}\selectfont \itshape #1}}
\newcommand{\Invited}{\textbf{[Invited]}}

% For fancy and multipage tables
\usepackage{tabularx}
\usepackage{ltablex}
\usepackage{environ}

% Macros to add links
\newcommand{\DOI}[1]{doi:\href{https://doi.org/#1}{#1}}
\newcommand{\Preprint}[1]{[preprint: \href{https://doi.org/#1}{doi.org/#1}]}
\newcommand{\Youtube}[1]{[recording: \href{https://youtu.be/#1}{youtu.be/#1}]}


\newcommand{\Duration}[2]{\textsuperscript{#1}/\textsubscript{#2}}
\newcommand{\Year}[1]{#1}
%\newcommand{\Ongoing}{Today}  
\newcommand{\Ongoing}{-}
\newcommand{\Future}{future}
\newcommand{\Review}{in review}
\newcommand{\Accepted}{accepted}
\newcommand{\OA}{{\bfseries [open access]}}


% Define command to insert month name and year as date
\usepackage{datetime}
\newdateformat{monthyear}{\monthname[\THEMONTH], \THEYEAR}


\setlength\parindent{0cm}

% Increase the line spacing
\renewcommand{\baselinestretch}{1.08}
% and the spacing between rows in tables

% Remove space between items in itemize and enumerate
\usepackage{enumitem}
\usepackage{multirow,bigdelim}
\usepackage{booktabs}
\usepackage{float}
\usepackage{array}
\setlist{nosep}

% Configure the font style for sections
\usepackage{xcolor}
\definecolor{awesome}{HTML}{96281B}
\definecolor{ieee}{HTML}{00629b}
\definecolor{acm}{HTML}{84cdf1}
\definecolor{springer}{HTML}{ee7d11}
\definecolor{elsevier}{HTML}{ff6c00}
\definecolor{wos}{HTML}{5e33bf}
\definecolor{scopus}{HTML}{ff6c00}
\definecolor{dblp}{HTML}{ffd700}
\definecolor{crossref}{HTML}{ef3340}
\definecolor{mdpi}{HTML}{3b5578}
% badger badger badger

\usepackage[most]{tcolorbox}


\newtcbox{\badge}[1][awesome]{
  on line, 
  arc=0pt,
  colback=#1!80!black,
  colframe=#1!80!black,
  fontupper=\scshape\small\color{white},
  boxrule=1pt, 
  boxsep=0pt,
  left=1pt,
  right=1pt,
  top=1pt,
  bottom=1pt
}

\newcommand{\wos}{\badge[wos]{WoS}}
\newcommand{\scopus}{\badge[scopus]{Scopus}}
\newcommand{\crossref}{\badge[crossref]{CrossRef}}
\newcommand{\dblp}{\badge[dblp]{DBLP}}

\newcommand{\ieee}{\badge[ieee]{IEEE}}
\newcommand{\acm}{\badge[acm]{ACM}}
\newcommand{\mdpi}{\badge[mdpi]{MDPI}}
\newcommand{\springer}{\badge[springer]{Springer}}
\newcommand{\elsevier}{\badge[elsevier]{Elsevier}}
% Set the spacing for sections
% HEADINGS
\usepackage{sectsty} 
\usepackage[normalem]{ulem} 
\sectionfont{\color{awesome}\mdseries\upshape\Large}
\subsectionfont{\color{awesome}\mdseries\scshape\normalsize} 
\subsubsectionfont{\color{awesome}\mdseries\upshape\large} 


% Set fancy headers
\usepackage{fancyhdr}
\pagestyle{fancy}
\fancyhf{}
\chead{
    \fontsize{8pt}{10pt}\selectfont
    \MyName
    \hspace{0.2cm}  \quad\adforn{24}\quad \hspace{0.2cm}
    \Title
    \hspace{0.2cm}  \quad\adforn{52}\quad  \hspace{0.2cm}
    \monthyear\today
}
\rhead{}
\usepackage{lastpage}
\cfoot{\fontsize{8pt}{0}\selectfont Page \thepage\ of \pageref*{LastPage}}
\renewcommand{\headrulewidth}{0pt}

\fancypagestyle{firststyle}
{
   \fancyhf{}
   \cfoot{\fontsize{8pt}{0}\selectfont Page \thepage\ of \pageref*{LastPage}}
}


% Metadata for the PDF output and control of hyperlinks
\usepackage[colorlinks=true]{hyperref}
\hypersetup{
    pdftitle={\MyName\ - \Title},
    pdfauthor={\MyName},
    linkcolor=blue,
    citecolor=blue,
    filecolor=black,
    urlcolor=awesome
}
%%%%%%%%%%%%%%%%%%%%%%%%%%%%%%%%%%%%%%%%%%%%%%%%%%%%%%%%%%%%%%%%%%%%%%%%%%%%%%%

\sloppy

\begin{document}

% No header for the first page
\thispagestyle{firststyle}

%%%%%%%%%%%%%%%%%%%%%%%%%%%%%%%%%%%%%%%%%%%%%%%%%%%%%%%%%%%%%%%%%%%%%%%%%%%%%%%
% HEADER
{\color{awesome}\fontsize{24pt}{0}\selectfont\MyName \hfill \adforn{74}}\\

\begin{minipage}[t]{0.595\textwidth}
	\Affiliation
	\\
	\Address
\end{minipage}
\begin{minipage}[t]{0.4\textwidth}
	\begin{flushright}
		Last updated: \monthyear\today
		\\
		ORCID: \href{https://orcid.org/\ORCID}{\ORCID}
		\\
		email: \href{mailto:\Email}{\Email}
		\\
		website: \href{https://\Website}{\Website}
	\end{flushright}
\end{minipage}
%\vspace{0.2cm}
\begin{center}
	\color{awesome} 
	\vspace{-0.85cm}
\end{center}

\paragraph{Summary} João Pedro Dias is part researcher on the thin line between hardware and software, and part Software Engineer. He has a BSc+MSc in Informatics and Computing Engineering by the Faculty of Engineering, University of Porto (FEUP). He earned his Ph.D. in Informatics Engineering from FEUP in 2022, while receiving a grant from the FCT. He maintains a Software Engineer position as a day-to-day job. Since 2017, he has been an Invited Assistant Professor at FEUP, where he teaches courses in Software Engineering, Operating Systems, among others. He has (co-)supervised +10 MSc dissertations and contributed to two public-funded projects at LIACC and INESC TEC. His research focuses on Internet-of-Things systems, software engineering, security and privacy, and his work has been published in several top-tier conferences and journals (h-index 18, i10-index 23). In his free time, he enjoys participating in Capture The Flag competitions, experimenting with Software-defined Radio, building web applications, reverse-engineering hardware, and photographing while wandering in nature.
\vspace{-1.1em}
%João Pedro Dias has a \textsc{BSc}+\textsc{MSc} in \textsc{Informatics and Computing Engineering} by the \textsc{Faculty of Engineering, University of Porto} (\textsc{FEUP}). He has a \textsc{Ph.D.} in \textsc{Informatics Engineering} by the same university circa 2022. He is an \textsc{Invited Assistant Professor} at \textsc{FEUP}, lecturing, since 2017, 8 different courses ranging from \textsc{Software Engineering} to \textsc{Operating Systems}. He has participate in 5$+$ MSc dissertations supervision. He previously worked as a researcher at \textsc{LIACC} and \textsc{INESC TEC} and he is now \textsc{Software Engineering Specialist} at \textsc{BUILT CoLAB}. He works in the area of Software Engineering, with a special interest in Design Patterns, Internet-of-Things, Security and Privacy, counting with 30$+$ published and indexed papers. In his leisure time, he can be found participating in Capture The Flag competitions, messing around with Software-defined Radio, doing web development, learning how to reverse engineer hardware and photographing while wandering in nature. 

%%%%%%%%%%%%%%%%%%%%%%%%%%%%%%%%%%%%%%%%%%%%%%%%%%%%%%%%%%%%%%%%%%%%%%%%%%%%%%%
\section*{Professional Appointments}

\years{2023---}
\textbf{Software Engineer}, Kuehne$+$Nagel\\
Porto IT Tech Hub, Porto, Portugal
\\
\years{2022---}
\textbf{Invited Assistant Professor}, Department of Informatics Engineering (DEI)
\newline
\FEUP
\\
\years{2021---2023}
\textbf{Software Engineering Specialist}, BUILT CoLAB \\
Collaborative Laboratory for the Built Environment of the Future, \Address{}
\\
\years{2017--2022}
\textbf{Invited Assistant Lecturer}, Department of Informatics Engineering (DEI)
\newline
\FEUP
\\
\years{2017--2021}
\textbf{Researcher}, Centre for Information Systems and Computer Graphics (CSIG)
\newline
Institute for Systems and Computer Engineering, Technology and Science (INESC TEC)
\newline
\Address{}
\\
\years{2016--2017}
\textbf{Researcher}, LIACC --- Artificial Intelligence and Computer Science Laboratory 
\newline
\FEUP
\\
\years{2015 (4 \footnotesize{mo.})} 
\textbf{Summer Engineering Intern}, ShiftForward, S.A. (\textit{a.k.a.} Velocidi, Inc., acquired by Kevel)\\
\Address{}
\\
\years{2014 (4 \footnotesize{mo.})}  
\textbf{IT Department Internship}, Centro Hospitalar do Tâmega e Sousa, E.P.E.
\\
Penafiel, \Address{}


%%%%%%%%%%%%%%%%%%%%%%%%%%%%%%%%%%%%%%%%%%%%%%%%%%%%%%%%%%%%%%%%%%%%%%%%%%%%%%%
\section*{Education}

\years{2016--2022}
\textbf{Doctoral Program in Informatics Engineering --- ProDEI (PhD)}
\newline \FEUP,  \textit{Unanimously Approved with Distinction (Cum Laude)}.
%\newline Supervision of Hugo Sereno Ferreira, PhD, and Co-Supervision of João Pascoal Faria, PhD 
\\
\years{2011--2016}  
\textbf{Integrated Masters in Informatics and Computing Engineering --- MIEIC (BSc + MSc)}
\newline \FEUP, \textit{GPA: 15 out of 20}.
\\
\years{2008--2011}
\textbf{Elementary School --- Sciences and Technologies}
\newline Externato de Vila Meã, \Address{}, \textit{GPA: 17 out of 20}.

%%%%%%%%%%%%%%%%%%%%%%%%%%%%%%%%%%%%%%%%%%%%%%%%%%%%%%%%%%%%%%%%%%%%%%%%%%%%%%%

\section*{Teaching}

\textbf{Integrated Master in Informatics and Computing Engineering (\textsc{MIEIC}) :: \textsc{FEUP} ::  5 years\footnote{MIEIC last edition was 20/21, then it became L.EIC + M.EIC.}}
\begin{table}[H]
	\renewcommand{\arraystretch}{0.25}
	\begin{tabular*}{\textwidth}{l @{\extracolsep{\fill}} ccccccc@{}}
		
		&                      &                      &  & \footnotesize{$^{17} / _{18}$} & \footnotesize{$^{18} / _{19}$} & \footnotesize{$^{19} / _{20}$}   & \footnotesize{$^{20} / _{21}$} \\ \midrule
		&                      &     Year                 & \small{ECTS}                & \multicolumn{4}{c}{Classes}         \\ \midrule \addlinespace[0.2cm]
		Object Oriented Programming Lab & \small{LPOO}                 & 2                    & 6                    & -                    & 1  & -   & -    \\\addlinespace[0.3cm]
		Operating Systems                      & \small{SOPE}                 & 2                    & 6                    & 2               & 1         & 2   & 2   \\\addlinespace[0.3cm]
		Software Engineering       & \small{ESOF}                 & 3                    & 6                  & -                    & -    & -          & 1 \\\addlinespace[0.3cm]
		Software Development Lab       & \small{LDSO}                 & 4                    & 7.5                  & -                    & 1    & -          & -   \\\addlinespace[0.3cm]
		Software Systems Architecture          & \small{ASSO}                 & 4                    & 6                    & -                    & 1    & -      & -   \\\addlinespace[0.2cm]\bottomrule
	\end{tabular*}
\end{table}

\textbf{Bachelor in Informatics and Computing Engineering (\textsc{L.EIC}) :: \textsc{FEUP} :: 3 years}
\begin{table}[H]
	\renewcommand{\arraystretch}{0.5}
	\begin{tabular*}{\textwidth}{l @{\extracolsep{\fill}} cccccc@{}}
		
		&                      &                      &  & \footnotesize{$^{21} / _{22}$} & \footnotesize{$^{22} / _{23}$} & \footnotesize{$^{23} / _{24}$}\\ \midrule
		&                      &     Year                 & \small{ECTS}                & \multicolumn{2}{c}{Classes}    &   -  \\ \midrule \addlinespace[0.2cm]
		Web Languages and Technologies & \small{LTW}                 & 2                    & 6                    & 2 & 2 & 2\\\addlinespace[0.3cm]
		Operating Systems                      & \small{SO}                 & 2                    & 6                    & 2 &   - & -  \\\addlinespace[0.2cm]\bottomrule
	\end{tabular*}
\end{table}

\textbf{Master in Informatics and Computing Engineering (\textsc{M.EIC}) :: \textsc{FEUP} :: 2 years}
\begin{table}[H]
	\renewcommand{\arraystretch}{0.5}
	\begin{tabular*}{\textwidth}{l @{\extracolsep{\fill}} ccccc@{}}
		
		&                      &                      &  & \footnotesize{$^{22} / _{23}$} & \footnotesize{$^{23} / _{24}$}\\ \midrule
		&                      &     Year                 & \small{ECTS}                & \multicolumn{1}{c}{Classes}   &      \\ \midrule \addlinespace[0.2cm]
		Large Scale Software Development                     & \small{DS}                 & 2                    & 6                    & 2    & 2   \\\addlinespace[0.2cm]\bottomrule
	\end{tabular*}
\end{table}

%%%%%%%%%%%%%%%%%%%%%%%%%%%%%%%%%%%%%%%%%%%%%%%%%%%%%%%%%%%%%%%%%%%%%%%%%%%%%%%
\section*{Publications}

Since 2017 I have (co-)authored \textbf{25 articles in conferences with scientific referee} and \textbf{7 articles in refereed periodics}. In total, these \textbf{31 publications} have attracted over \textbf{633 citations}, resulting in an \textbf{h-index of 17} and \textbf{i10-index of 23}\footnote{Citations and indexes were obtained by using Google Scholar, \textsc{id: \href{https://scholar.google.com/citations?user=sQ2vKI0AAAAJ}{sQ2vKI0AAAAJ}}, on 31 July, 2023}. Indexation of the articles in Scopus, Web Of Science (WoS), DBLP, and CrossRef databases is represented using badges.

\subsection*{Peer-Reviewed Journals}

\begin{etaremune}
	\renewcommand\labelenumi{\footnotesize\theenumi .}
    \item \fullcite{RIBEIRO2024112081} \scopus \wos \dblp \crossref
	\item \fullcite{Sosym2022} \scopus \wos \dblp \crossref
	\item \fullcite{DiasIoT2022} \scopus \wos \dblp \crossref
	\item \fullcite{AccessReis2021} \scopus \wos \dblp \crossref
	\item \fullcite{lago2020journal} \scopus \wos \dblp \crossref
	\item \fullcite{ccis19a} \scopus \wos \dblp \crossref
	\item \fullcite{ccis19b} \scopus \wos \dblp \crossref
	\item \fullcite{jias18} \wos \dblp
\end{etaremune}

\subsection*{Peer-Reviewed Conference Proceedings}

\begin{etaremune}
	\renewcommand\labelenumi{\footnotesize\theenumi .}
    \item \fullcite{verdi2023}
	\item \fullcite{icse2022} \scopus \wos \dblp 
	\item \fullcite{icse2021} \scopus \wos \dblp \crossref
	\item \fullcite{DannyUserStudy21} \dblp \crossref
	\item \fullcite{MargaridaSurvey21} \dblp \crossref
	\item \fullcite{Margarida2020} \wos \dblp \crossref
	\item \fullcite{europlop2020} \scopus \dblp \crossref
	\item \fullcite{DiogoTorres2020} \scopus \wos \dblp \crossref
	\item \fullcite{Piedade2020} \scopus \dblp \crossref
	\item \fullcite{selfheal20} \scopus \dblp \crossref
	\item \fullcite{Lago2020} \scopus \dblp \crossref
	\item \fullcite{europlop19} \scopus \wos \dblp \crossref
	\item \fullcite{Aguiar2019} \scopus \dblp \crossref
	\item \fullcite{enasePedro19} \scopus \wos \dblp \crossref
	\item \fullcite{enaseGil19} \scopus \wos \dblp \crossref
	\item \fullcite{Pinto18} \scopus \wos \dblp \crossref
	\item \fullcite{DiasICST18} \scopus \wos \dblp \crossref
	\item \fullcite{SEDES2018} \scopus \wos \dblp \crossref
	\item \fullcite{ISC22018} \scopus \dblp \crossref
	\item \fullcite{IAS2018health} \scopus \dblp \crossref
	\item \fullcite{pinto2018blockchainbased} \scopus \dblp \crossref
	\item \fullcite{Ramadas2017}
	\item \fullcite{Dias2017} \scopus \dblp \crossref
	\item \fullcite{Duarte2017} \scopus \dblp \crossref
	\item \fullcite{Diasecommerce2017} \scopus \wos \dblp \crossref
\end{etaremune}

\subsection*{Informal Publications \textit{\&} Others}

\begin{etaremune}
	\renewcommand\labelenumi{\footnotesize\theenumi .}
	\item \fullcite{xp2020survey}
	\item \fullcite{ijup18}
	\item \fullcite{ijup17}
	\item \fullcite{ijup16}
\end{etaremune}

\subsection*{Preprints}

\begin{etaremune}
	\renewcommand\labelenumi{\footnotesize\theenumi .}
	\item  \fullcite{DBLP:journals/corr/abs-2009-05828}
\end{etaremune}


\subsection*{Dissertations}


\years{2022} \fullcite{feupphd}
\\[0.42em]
\years{2016} \fullcite{feupmsc}



%%%%%%%%%%%%%%%%%%%%%%%%%%%%%%%%%%%%%%%%%%%%%%%%%%%%%%%%%%%%%%%%%%%%%%%%%%%%%%%

\section*{Educational Activities}

Since 2016, I have been (co-)supervisor of +\textbf{10 M.Sc. dissertations} and of \textbf{6 B.Sc. capstone projects}. Most of my (co-)supervisions consist of students of the Informatics and Computation Engineering degree by \FEUP, with several collaborations with colleagues as main or co-supervisor. Most of these supervisions are pursued in an academic context (\textit{i.e.}, within the institution that confers the degree). Others took place at companies/external institutions, including: Critical Manufacturing, Smartex.ai, DevScope, Kuehne+Nagel, and BUILT CoLAB.

\subsection*{Master's Dissertation Supervisions}

\begin{etaremune}
	\renewcommand\labelenumi{\footnotesize\theenumi .}
    \item \textit{Incremental Load Test on Event-Driven Architectures with Structured Data} by Ângela Coelho (2023--2024).
	Co-supervisor: António Sernadas, MSc (Kuehne-Nagel). \FEUP.
    \item \textit{AsyncAPI-First Design for Event-Driven Architectures: Improve Developer Experience} by José Silva (2023--2024).
	Co-supervisor: António Sernadas, MSc (Kuehne-Nagel). \FEUP.
    \item \textit{Adopting Continuous Deployment of WebAssembly Applications in Internet-of-Things Systems} by Joel Fernandes (2023--2024).
	Co-supervisor: André Restivo, PhD. \FEUP.
	\item \textit{Leveraging Context-awareness in IoT Systems: An Approach with Voice Assistants} by Daniel Gonçalves (2022--2023).
	Co-supervisor: André Restivo, PhD. \FEUP.
	
	\item \textit{Adopting Containers in Microcontrollers for the IoT} by Xavier Pisco (2022--2023).
	Co-supervisor: André Restivo, PhD. \FEUP.

	\item \textit{Privacy-aware Rule Conflict Prediction in Home Automation} by Pedro Ponte (2022--2023).
	Main-supervisor: André Restivo, PhD. \FEUP.

	\item \textit{Micro-containerization in Microcontrollers for the IoT} by Eduardo Ribeiro (2021--2022). \textit{handle}: \href{http://hdl.handle.net/10216/142728}{10216/142728}.
	Co-supervisor: Hugo Sereno Ferreira, PhD. \FEUP.
	
	\item \textit{MQTT Chaos Engineering for Self-Healing IoT Systems} by Miguel Duarte (2020--2021), \textit{ref.}: \href{https://hdl.handle.net/10216/135470}{10216/135470}. Main supervisor: André Restivo, PhD. \FEUP.
	
	\item \textit{Visual Programming Language for Orchestration with Docker} by Bruno Piedade (2019--2020), \textit{ref.}: \href{http://hdl.handle.net/10216/128974}{10216/128974}. Main supervisor: Filipe Correia, PhD. \FEUP.
	
	\item \textit{Implementing a Multi-Approach Debugging of Industrial IoT Workflows} by Andreia Rodrigues (2018--2019), \textit{ref.}: \href{http://hdl.handle.net/10216/122773}{10216/122773}. Main supervisor: Hugo Sereno Ferreira, PhD; Collaboration with \href{https://www.criticalmanufacturing.com/}{Critical Manufacturing}. \FEUP.
	
	\item \textit{Dynamic Allocation of Serverless Functions in IoT Environments} by Duarte Pinto (2017--2018), \textit{ref.}: \href{http://hdl.handle.net/10216/114149}{10216/114149}. Main supervisor: Hugo Sereno Ferreira, PhD. \FEUP.
	
	\item \textit{Blockchain as a PKI for Ownership Control of IoT Devices} by Guilherme Pinto (2017--2018), \textit{ref.}: \href{http://hdl.handle.net/10216/114144}{10216/114144}. Main supervisor: Hugo Sereno Ferreira, PhD. \FEUP.
	
	\item \textit{Interoperability In Software Applications For Smart Cities: Towards A Reference Architecture} by José Pedro Pinto (2016--2017), \textit{ref.}:  \href{http://hdl.handle.net/10216/107529}{10216/107529}. Main supervisor: Rosaldo Rossetti, PhD. \FEUP.
\end{etaremune}

\subsection*{Bachelor's Capstone Project Supervisions}

\begin{etaremune}
	\renewcommand\labelenumi{\footnotesize\theenumi .}
 \item \textit{IoT-based Real-Time Helpdesk Occupation Status V2.0}. Co-supervisor: João Paiva Cardoso. André Costa, Fábio Sá, Lourenço Gonçalves, Marcos Pinto. \textit{Group Project} (2022--2023). \FEUP.
	\item \textit{Hardware Simulation to Assist Software Development for the Textile Industry}. Sérgio Estêvão, \textit{Individual Internship} at \href{https://www.smartex.ai/}{Smartex.ai} (2021--2022). \FEUP. 
	\item \textit{Research Project Management as Scrum Master in the DevScope Academy}. João Baltazar, \textit{Individual Internship} at \href{https://www.devscope.net/}{DevScope} (2021--2022). \FEUP. 
	
	\item \textit{Multi-purpose License Manager}. Daniel Gonçalves, Isla Cassamo, André Malheiro, \textit{Group Project} in collaboration with \href{https://www.builtcolab.pt/}{BUILT CoLAB} (2021--2022). \FEUP.
	
	\item \textit{IoT-based Real-Time Helpdesk Occupation Status}. Main supervisor: João Paiva Cardoso, PhD. João Pinho, Pedro Santos, Tiago Gonçalves, Diogo Vale, \textit{Group Project} (2021--2022). \FEUP.
	
	\item \textit{IoT-enhanced HMI --- Embedded Mouse}. Main supervisor: Pedro Diniz, PhD. Bruno Mendes, José Costa, David Preda, Fernando Rego, \textit{Group Project} (2021--2022). \FEUP.
	
	\item \textit{Ubiquitous Sensing --- Sensing, Storing and Processing Plant Sensor Data}. Main supervisor: Pedro Diniz, PhD. Nuno Alves, Miguel Amorim, Mário Travassos, Ricardo Ferreira, \textit{Group Project} (2021--2022). \FEUP.
\end{etaremune}

%%%%%%%%%%%%%%%%%%%%%%%%%%%%%%%%%%%%%%%%%%%%%%%%%%%%%%%%%%%%%%%%%%%%%%%%%%%%%%%
\section*{Invited Talks, Presentations \textit{\&} Workshops}

Slides and abstracts for most of the presentations are available at
\href{https://jpdias.me/talks}{jpdias.me/talks}.
\\

\years{2024}  
\fullcite{devopsKN24}
\\[0.42em]
\years{2023}  
\fullcite{apidays23}
\\[0.42em]
\fullcite{ptpc23}
\\[0.42em] 
\fullcite{isep23}
\\[0.42em]
\years{2021}  
\fullcite{sinf2021}
\\[0.42em] 
\fullcite{update2021}
\\[0.42em]
\years{2020}  
\fullcite{codeweek20}
\\[0.42em]
\years{2019}  
\fullcite{oposec19}
\\[0.42em]
\fullcite{pixelscamp3}
\\[0.42em]
\years{2017}  
\fullcite{pixelscamp2}
\\[0.42em]
\years{2016}  
%\Invited{}
\fullcite{oposec16}
\\[0.42em]
\years{2015}  
%\Invited{}
\fullcite{feup15}


%%%%%%%%%%%%%%%%%%%%%%%%%%%%%%%%%%%%%%%%%%%%%%%%%%%%%%%%%%%%%%%%%%%%%%%%%%%%%%%

\section*{Academic Service \textit{\&} Affiliations}


Since 2017, I have participated as a reviewer\footnote{All reviews are verified on Web of Science: \href{https://www.webofscience.com/wos/author/record/O-4128-2018}{wos/O-4128-2018}.} in \textbf{7 periodics} and \textbf{7 international
	conferences}. I also have been member of \textbf{11 committees} of international and national conferences, and have been a volunteer at 3 different scientific events.

\subsection*{Reviewer for International Conferences}

\begin{itemize}[leftmargin=0pt,labelsep=10pt,itemsep=1.4pt,parsep=2pt]
\small
    \item \ieee{} International Smart Cities Conference 2021 (IEEE ISC2 2021) \small{(1$\times$ Full Paper)}
	\item 2\textsuperscript{nd} International Conference on Artificial Intelligence, Information Processing and Cloud Computing (AIIPCC 2021) \small{(1$\times$ Full Paper)}
	\item 4\textsuperscript{th} International Conference on Computer Science and Application Engineering (CSAE 2020) \small{(1$\times$ Full Paper)}
	\item \ieee{} International Symposium on Personal, Indoor and Mobile Radio Communications (IEEE PIMRC 2019) \small{(1$\times$ Full Paper)}
	\item \ieee{} Intelligent Vehicles Symposium (IEEE IV 2017) \small{(1$\times$ Full Paper)}
	\item \ieee{} Summer School on Smart Cities (IEEE S3C 2017) \small{(2$\times$ Full Papers, 1$\times$ Short Paper)}
	\item \ieee{} 20\textsuperscript{th} International Conference on Intelligent Transportation Systems (IEEE ITSC 2017) \small{(2$\times$ Full Paper)}
\end{itemize}

\subsection*{Reviewer for International Scientific Journals}

\begin{itemize}[leftmargin=0pt,labelsep=10pt,itemsep=1.4pt,parsep=2pt]
\small{
    \item \acm{} Proceedings of the ACM on Interactive, Mobile, Wearable and Ubiquitous Technologies (1$\times$)
    \item \ieee{} Access (1$\times$)
    \item \ieee{} Intenert of Things Journal (1$\times$)
    \item \elsevier{} Journal of Information Security and Applications (4$\times$)
    \item \springer{} Soft Computing (1$\times$)
	\item \springer{} Journal of Grid Computing (1$\times$)
    \item \badge{Taylor and Francis} Journal of Experimental \& Theoretical Artificial Intelligence (1$\times$)
    \item \badge{Taylor and Francis} International Journal of Human-Computer Interaction (1$\times$)
	\item \mdpi{} Journal of Cybersecurity and Privacy (1$\times$)
    \item \mdpi{} Electronics (1$\times$)
 	\item \mdpi{} Symmetry (1$\times$)
	\item \badge[blue]{Oxford} Interacting with Computers: Interdisciplinary Journal of Human-Computer Interaction (1$\times$)
	\item \badge[mdpi]{ITB} Journal of ICT Research and Applications (1$\times$)
	}
\end{itemize}

\subsection*{Member of Committees}

\years{2023} 
\textbf{Program Committee Member}, ICSOFT --- International Conference on Software Technologies, 18\textsuperscript{th} Edition\\[0.42em]
\textbf{Organizing Committee}, 5\textsuperscript{th} International Workshop on Software Engineering Research \textit{\&} Practices for the Internet of Things (SERP4IoT 2023), co-located with the 44\textsuperscript{th} ACM/IEEE International Conference on Software Engineering (ICSE 2023)
\\[0.42em]
\years{2022} \textbf{Publicity Chair}, 4\textsuperscript{th} International Workshop on Software Engineering Research \textit{\&} Practices for the Internet of Things (SERP4IoT 2022), co-located with the 44\textsuperscript{th} ACM/IEEE International Conference on Software Engineering (ICSE 2022)
\\[0.42em]
\years{2019} \textbf{Track Co-Chair}, Symposium on Informatics Engineering, part of the 3\textsuperscript{rd} Doctoral Congress Engineering (DCE19), \FEUP
\\[0.42em]
\textbf{Program Committee Member}, Talk a Bit --- Student Organized Tech Conference, 7\textsuperscript{th} Edition
\\[0.42em]
\years{2018} \textbf{Makerspace Chair}, 19\textsuperscript{th} International Conference on Agile Software Development (XP 2018)
\\[0.42em] 
\textbf{Conference Organization Chair}, 13\textsuperscript{th} Doctoral Symposium in Informatics Engineering (DSIE'18)
\\[0.42em] 
\textbf{Advisor Board Member}, Talk a Bit --- Student Organized Tech Conference, 6\textsuperscript{th} Edition
\\[0.42em]
\years{2017} \textbf{Program Committee Member}, 2017 IEEE 1\textsuperscript{st} Summer School on Smart Cities (IEEE S3C 2017)
\\[0.42em]
\textbf{Web Communication Co-chai}r, The Second ACM International Workshop on Smart, Autonomous, and Connected Vehicular Systems and Services (ACM CarSys 2017), part of the The 23\textsuperscript{rd} Annual International Conference on Mobile Computing and Networking (MobiCom 2017)
\\[0.42em]
\textbf{Advisor Board Member}, Talk a Bit --- Student Organized Tech Conference, 5\textsuperscript{th} Edition (2017)
\\[0.42em]
\years{2016} \textbf{Conference Chair}, Talk a Bit --- Student Organized Tech Conference, 4\textsuperscript{th} Edition (2016)
\\[0.42em]
\textbf{Web Communication Co-chair}, 2016 IEEE 19\textsuperscript{th} International Conference on Intelligent Transportation Systems (ITSC 2016)


\subsection*{Volunteer Activities}


\years{2017} 
\textbf{Student Volunteer (Media Coverage)}, 18\textsuperscript{th} EPIA Conference on Artificial Intelligence (EPIA 2017), Porto, Portugal
\\[0.42em]
\years{2014} 
\textbf{Student Volunteer}, Inter-University Programming Marathon (MIUP 2014), \Address{}
\\[0.42em]
\textbf{Student Volunteer}, INForum - Symposion on Informatics (INForum 2014), \Address{}


\subsection*{Affiliations}


\years{2017--} \textbf{Software Engineering Group} at the Department of Informatics Engineering, Faculty of Engineering, University of Porto.
\\[0.42em]
\years{2017--} \textbf{Institute for Systems and Technologies of Information, Control and Communication} (INSTICC).
\\[0.42em]
\years{2019--2020} \textbf{Association for Computing Machinery} (ACM), Student Membership.
\\[0.42em]
\years{2013--2018}  \textbf{Institute of Electrical and Electronics Engineers} (IEEE), Student Membership as part of the University of Porto Student Branch (IEEE UP SB).


%%%%%%%%%%%%%%%%%%%%%%%%%%%%%%%%%%%%%%%%%%%%%%%%%%%%%%%%%%%%%%%%%%%%%%%%%%%%%%%
\section*{Grants, Honors \textit{\&} Awards}

\subsection*{Grants}

%\textbf{PhD Scholarship}\marginpar[\Duration{2019}{2021}]{} , by FCT --- Fundação para a Ciência e a Tecnologia (Portuguese Foundation for Science and Technology), \textit{ref.} SFRH/BD/144612/2019, 2 years. 
\years{2019--2021} 
\textbf{PhD Scholarship}, by FCT --- Fundação para a Ciência e a Tecnologia (Portuguese Foundation for Science and Technology), \textit{ref.} SFRH/BD/144612/2019, 2 years.
\\[0.42em]
\years{2017--2018} 
\textbf{NanoSTIMA --- Macro-to-Nano Human Sensing: Towards Integrated Multimodal Health Monitoring and Analytics} at INESC TEC, \textit{ref.} NORTE-01-0145-FEDER-000016, 12 months, \href{http://nanostima.inesctec.pt/}{nanostima.inesctec.pt}
\\[0.42em]
\years{2016--2017} 
\textbf{PEDDIR DEMO: Weigth in Motion and Wheels Defect Detection (Sistema de Pesagem Dinâmica e Deteção de Irregularidades dos Rodados)} at LIACC, \textit{ref.} NORTE-01-0247-FEDER-006397, 6 months, \href{http://peddir.eu/}{peddir.eu}



\subsection*{International Honors \textit{\&} Awards}


\years{2017}  
\textbf{427\textsuperscript{th} \small{(of 3350)}}, IEEEXtreme --- 24h Programming Competition 11.0, 
\href{http://sites.ieee.org/xtreme/}{ieee.org/xtreme} (Online)
\\[0.42em]
\textbf{847\textsuperscript{th} \small{(of 2815)}}, Google Hash Code 2017, 
\href{http://hashcode.withgoogle.com}{hashcode.withgoogle.com} (Online)
\\[0.42em]
\years{2015}  
\textbf{189\textsuperscript{th} \small{(of 2477)}}, IEEEXtreme --- 24h Programming Competition 9.0, 
\href{http://sites.ieee.org/xtreme/}{ieee.org/xtreme} (Online)
\\[0.42em]
\years{2014} 
\textbf{Honorable Mention}, IEEEmadC --- IEEE Mobile Application Dev. Contest, \href{http://ieeemadc.org}{ieeemadc.org} (Online)
\\[0.42em]
\years{2013}  
\textbf{126\textsuperscript{th} \small{(of 2570)}}, IEEEXtreme --- 24h Programming Competition 7.0, 
\href{http://sites.ieee.org/xtreme/}{ieee.org/xtreme} (Online)
\\[0.42em]
\textbf{2\textsuperscript{nd} Place} at European Final (1\textsuperscript{st} at Local and National rounds), European BEST Engineering Competition, \href{http://ebec.best.eu.org}{ebec.best.eu.org}, Warsaw, Poland.


\subsection*{Domestic Honors \textit{\&} Awards}
\years{2022} 
\textbf{11\textsuperscript{th} \small{(of 109 Participants)}},  Portuguese Cybersecurity Competition by hackrocks, \href{https://cyberpt.hackrocks.com/}{cyberpt.hackrocks.com}. Online (Portugal).
\\[0.42em]
\years{2021} 
\textbf{7\textsuperscript{th} \small{(of 372 Teams)}}, Portuguese Stage of the 2021 MetaRed CTF Championship by RCTS-FCCN (FCT Portugal) and Metared, \href{https://cert.rcts.pt/}{cert.rcts.pt}. Online (Portugal).
\\[0.42em]
\years{2020}  
\textbf{6\textsuperscript{th} \small{(of 51 Teams)}}, CIBER PERSEU 2020 Portguese Army Exercise --- Capture The Flag component by \href{https://www.ihacklabs.com/}{iHackLabs} \textit{\&} \href{http://exercito.pt/}{exercito.pt}. Online (Portugal).
\\[0.42em]
\textbf{9\textsuperscript{th} \small{(of 209 Teams)}}, Portuguese Stage of the 2020 MetaRed CTF Championship by RCTS-FCCN (FCT Portugal) and Metared, \href{https://cert.rcts.pt/}{cert.rcts.pt}. Online (Portugal).
\\[0.42em]
\textbf{3\textsuperscript{th} \small{(of 6 Teams)}}, Capture The Flag by BSidesPorto, \href{https://bsidesporto.com/}{bsidesporto.com}. Online (Portugal).
\\[0.42em]
\years{2018}  
\textbf{6\textsuperscript{th} \small{(of 31 Participants)}}, CIBER PERSEU 2018 Portguese Army Exercise --- Capture The Flag component by \href{https://www.diateam.net}{\texttt{diateam}} \textit{\&} \href{http://exercito.pt/}{exercito.pt}. Portuguese Military Academy, Lisbon, Portugal.
\\[0.42em]
\years{2011} 
\textbf{1\textsuperscript{st} Place}, PmatE Mat12 --- Mathematics Competition, \href{http://pmate.ua.pt/}{pmate.ua.pt}. University of Aveiro, Portugal.
\\[0.42em]
\years{2009}  
\textbf{2\textsuperscript{nd} Place}, PmatE Fis12 --- Physics Competition, 
\href{http://pmate.ua.pt/}{pmate.ua.pt}. University of Aveiro, Portugal.


%%%%%%%%%%%%%%%%%%%%%%%%%%%%%%%%%%%%%%%%%%%%%%%%%%%%%%%%%%%%%%%%%%%%%%%%%%%%%%%

\section*{Additional Training}

\years{2024}  
\textbf{Terraform Associate (003)} by HashiCorp Certified (Online with Proctored Exam)
\\[0.42em]
\years{2022}  
\textbf{Redis Certified Developer}, with completion of the modules \textit{Introduction to Redis Data Structures} (RU101), \textit{Redis Streams} (RU202), \textit{Redis Security} (RU330), by Redis University (Online)
\\[0.42em]
\years{2020}  
\textbf{Introduction to Reverse Engineering with Ghidra} by Matthew Alt of Voidstar Security, 4-hours training, Hackaday-U (Online)
\\[0.42em]
\textbf{Autopsy Basics and Hands On} by Brian Carrier, 8-hours training, Basis Technology Corp. (Online)
\\[0.42em]
\textbf{Practical Ethical Hacking - The Complete Course} by Heath Adams, 24.5-hours training, TCM Security, Inc., Udemy (Online)
\\[0.42em]
\textbf{Python for Penetration Testers} by Cristi Zot, 2-hours training, Udemy (Online)
\\[0.42em]
\years{2018} 
\textbf{1\textsuperscript{st} UPTEC School on the Future of Computing}, 1-week training, UPTEC – Science and Technology Park of University of Porto
\\[0.42em]
\years{2017} 
\textbf{CyberTalk: Security - Business Enablement}, 4-hours training, Multicert – Electronic Certification Services, S.A.
\\[0.42em]
\textbf{Project Management}, 12-hours training, ERASMUS+ Mobility for teaching at \FEUP
\\[0.42em]
\years{2016} 
\textbf{AWS Technical Essentials Online}, 1-day training, Amazon Web Services – Training and Certification
\\[0.42em]
\years{2015} 
\textbf{Practical Workshop in Agile Project Management}, 4-hours training, Strongstep - Innovation in Software Quality
\\[0.42em]
\textbf{CERN Spring Campus}, 1-week training, CERN - European Organization for Nuclear Research
\\[0.42em]
\years{2014} 
\textbf{M101JS: MongoDB for Node.js Developers}, 25-hours training, MongoDB University (Online)



\section*{Prototypes}

The prototypes are released under open-source licences, most of them being available through my Github profile:
\href{https://github.com/jpdias/}{github.com/jpdias}.
\\

\years{2022---} 
\textbf{Stateless Home Automation Dashboard Generation Tool}, 
\href{https://github.com/BUILT-CoLAB/mqttdash}{github.com/BUILT-CoLAB/mqttdash}
\newline
A simple codegen tool for developing single-page dashboards using MQTT-over-Websockets. 
\\[0.42em]
\textbf{Software License Manager and Verification API}, 
\href{https://github.com/BUILT-CoLAB/license-key-manager}{github.com/BUILT-CoLAB/license-key-manager}
\newline
A Flask-based (Python) platform to manage software licenses with an API for license verification.
\\[0.42em]
\years{2019---} 
\textbf{Self-Healing Extensions for Node-RED}, 
\href{https://github.com/jpdias/node-red-contrib-self-healing}{github.com/jpdias/node-red-contrib-self-healing}
\newline
A collection of self-healing nodes to the Node-RED development environment.
\\[0.42em]
\years{2016--2017} 
\textbf{Botnet-Lab}, 
\href{https://github.com/jpdias/botnet-lab}{github.com/jpdias/botnet-lab}
\newline
An IRC based tool for testing the capabilities of a botnet written in Python 2.7. 
\\
\years{2016} 
\textbf{Portfolio-Photo}, 
\href{https://github.com/jpdias/portfolio-photo}{github.com/jpdias/portfolio-photo}
\newline
An YAML and Jekyll (Ruby) dynamically generated photography portfolio using the Flickr API.
\\[0.42em]
\years{2015} 
\textbf{Deploy Logger Service (Codename Malamute)}, 
\href{https://github.com/ShiftForward/malamute}{github.com/ShiftForward/malamute}
\newline
Application and integration API built for documenting and record the release of new versions of projects (deploys) and its modules. Developed during the summer internship at ShiftForward.
\\[0.42em]
\years{2014--2015} 
\textbf{Turing Machine}, 
\href{https://github.com/jpdias/TuringMachine/}{github.com/jpdias/TuringMachine}
\newline
A web-based Turing Machine with a pre-defined language, parser, and interpretation capabilities.

%%%%%%%%%%%%%%%%%%%%%%%%%%%%%%%%%%%%%%%%%%%%%%%%%%%%%%%%%%%%%%%%%%%%%%%%%%%%%%%


\section*{Pro-Bono \textit{\&} Other Activities}

\subsection*{Academic}


\years{2013--2015} 
\textbf{Executive Committee Member (Secretary)}, IEEE University of Porto Student Branch
\newline
Organization of several technical workshops and other events including \emph{IEEE Day}, \emph{IEEEXtreme} and \emph{IEEEmadC}.


\subsection*{Humanitarian}


\years{2013--2015} 
\textbf{Founder Member}, \emph{Luz da Vida} Youth Group
\newline
Organization of activities on behalf of the village community (Travanca, Amarante, Portugal).

\subsection*{Hobbies}


\years{2016--2017} 
\textbf{Football Player}, Grupo Desportivo e Recreativo de Travanca 
\newline
F.A.D.A. Amateur League, Amarante, \Address{}
\\[0.42em]
\years{2008---} 
\textbf{Amateur Photographer} --- Wildlife, Landscapes and Travel
\newline
- \textit{Winner of the quest \textit{Landscapes from Above}}, promoted by 500px, \href{https://web.500px.com/quests/1000000106/landscapes-from-above}{500px.com/quests/1000000106}, 2020
\newline
- \textit{One of the 50 best pictures of \#MySocialEurope photo competition - How should the Future of social Europe look like?}, promoted by the European Commission, 2019
\newline
- \textit{Finalist at the Norte Multimédia Contest - Photography Category}, promoted by the Norte Portugal Regional Coordination and Development Commission (CCDR-N), exposed at Trindade Metro Station (Porto), 2018
\newline
- \textit{Finalist at the PORTUGAL NUM ISQUEIRO BIC Contest}, promoted by the I Love My BIC initiative by BIC Portugal, 2018
\newline
- \textit{Exposition at Movimento d'Artes}, promoted by the Cooperativa Cultural Vinha das Artes, 2010

%%%%%%%%%%%%%%%%%%%%%%%%%%%%%%%%%%%%%%%%%%%%%%%%%%%%%%%%%%%%%%%%%%%%%%%%%%%%%%%

\section*{Languages}

\TablePad
\begin{tabularx}{\textwidth}{@{}p{0.15\textwidth} p{0.85\textwidth}@{}}
	Portuguese & Native                    
	\\
	English    & B Level (Self-Evaluation) 
\end{tabularx}

%\begin{center}
%	\color{awesome} \Huge \adforn{17}
%	\vspace{-0.5cm}
%\end{center}

\end{document}